% =====================================================================
% Section 2 -- Model. Written in step 6.3 from model.tex (phase 2) with
% the corrections of notes 4.9 points 4, 9, 11 (gravity only in the
% static state), 12 and 19 (notation of step 6.2: delta, m', lambda,
% delta_c). Points 1-3 and 6 are applied in Section 3.
% Numbers quoted from the literature (lodewijks2013, Fig. 8) are typed;
% computed numbers come from paper_numbers.tex.
% =====================================================================
\section{Model}\label{sec:model}

% =====================================================================
% notation.tex -- table of notation (included in Section 2).
% Draft of step 6.2 (completed for Sections 2-3 in step 6.3), from model.tex and the notes (sections 4.1,
% 4.14-4.26). Complete in 6.3-6.6 as symbols enter the text; prune in 6.8.
%
% Decisions of step 6.2 (clashes resolved):
%   - Periods are t (s) and tau (units of L/c_r), like t_a and tau_a:
%     t_1, tau_1 for the loop, t_s, tau_s for a fixed-free strand,
%     t_{B1}, t_{A1} for modes followed by strand. T stays for tensions
%     (T_0, T_d, T_t, and T_1, T_2 at the drive as in DIN 22101).
%   - Transit time of strand B: tau_B = 1 - xi + gamma (was t_B).
%   - Drive factor written e^{mu theta}; E is only Young's modulus.
%     Wrap angle theta (DIN 22101 uses alpha, taken here by the material
%     coupling factor); friction coefficient mu without subscript, while
%     line densities always carry one (mu_r, mu_c).
%   - Inclination of the belt path delta(s), as in DIN 22101 (model.tex
%     used theta(s)); carriage inclination delta_c (model.tex used kappa
%     for its sine, which clashed with the wavenumbers kappa_j).
%   - Weighing line density m' (DIN 22101); m is kept for masses
%     (modal mass m_k, strand masses m_A, m_B, drive mass ratio m_d).
%   - Rigging ratio lambda (model.tex used i, the imaginary unit).
%   - T_t/F_U gets no symbol (the notes of step 5.4 used tau).
% =====================================================================
\begin{table}[p]
\centering
\caption{Notation. Dimensionless quantities are scaled with the strand length $L$,
the return-strand transit time $L/c_r$, the peak drive acceleration $a_m$ and the
inertial force $\mu_rLa_m$.}
\label{tab:notation}
\footnotesize
\begin{tabular}{@{}l>{\raggedright\arraybackslash}p{0.72\linewidth}@{}}
\toprule
Symbol & Meaning\\
\midrule
\multicolumn{2}{@{}l}{\emph{Geometry}}\\
$L$ & path length of each strand (carry and return)\\
$\sigma$ & geometric coordinate from the head pulley in the direction of travel; return strand $0\le\sigma<L$\\
$\sigma_d$, $\sigma_t$ & positions of the drive and take-up pulleys\\
$s$, $x=s/L$ & loop coordinate from the drive exit; $s=0$ and $s=2L$ are the faces of the drive pulley\\
$s_t$, $\xi=s_t/L$ & take-up position from the drive exit\\
$\ell_1$, $\ell_2$ & return belt between head and drive, and between take-up and tail (units of $L$)\\
$\delta$ & local inclination of the belt path ($\delta>0$ uphill)\\
\midrule
\multicolumn{2}{@{}l}{\emph{Belt, drive and take-up}}\\
$EA$ & axial stiffness of the belt ($E$, Young's modulus)\\
$\mu_r$, $\mu_c$ & inertial line densities (belt, reduced idler mass, coupled material)\\
$\mw_r$, $\mw_c$ & line densities that contribute to weight (belt and material)\\
$c_r$, $c_c$ & wave speeds, $c_j=\sqrt{EA/\mu_j}$; $Z_r=\mu_rc_r$, impedance of the return strand\\
$\alpha$ & material coupling factor \citep{lodewijks2002}\\
$t_v$ & Kelvin--Voigt retardation time\\
$\Mtu$; $M_w$, $M_c$ & take-up mass referred to the carriage; counterweight and carriage masses\\
$\nstr$ & belt strands that carry the take-up pulley\\
$\lambda$ & rigging ratio (counterweight travel per unit carriage travel; $\lambda=1$ when hung directly)\\
$\delta_c$ & inclination of the take-up carriage\\
$\Tt$ & static belt tension at the take-up\\
$T_1$, $T_2$ & belt tensions at the drive entry and exit (tight and slack side)\\
$\Wr=g\mu_rL$ & weight of the return strand\\
$F_U$ & running peripheral force at the drive (DIN 22101)\\
$\egrip$ & drive factor (friction coefficient $\mu$ between belt and pulley, wrap angle $\theta$)\\
$m_d$, $\hat c$ & reduced drive mass over $\mu_rL$ and slip damping over $\mu_rc_r$ (drive without speed control)\\
\midrule
\multicolumn{2}{@{}l}{\emph{Response}}\\
$w$, $U$ & elastic displacement of the belt; rigid motion imposed by the drive\\
$y$, $y_c$ & take-up displacement (two-strand equivalent, positive when the loop lengthens); carriage displacement\\
$T_0$, $T_d$ & static and dynamic belt tension; $\hat T=T_d/(\mu_rLa_m)$\\
$V$, $a$, $a_m$ & drive velocity, acceleration and peak acceleration\\
$t_a$, $\tau_a$ & start (acceleration) time; $\tau_a=c_rt_a/L$\\
$r$, $\varphi$ & motion resistance per unit length; its onset function\\
\midrule
\multicolumn{2}{@{}l}{\emph{Dimensionless parameters and modal quantities}}\\
$\gamma=c_r/c_c$ & wave-speed ratio, $\gamma^2=\mu_c/\mu_r$\\
$\beta$ & take-up mass ratio, $4\Mtu/(\nstr^2\mu_rL)=(4/\nstr)\lambda\,\Tt/\Wr$ (carriage mass neglected)\\
$\hat\zeta$ & damping parameter, $c_rt_v/(2L)$; modal damping $\zeta_k=\hat\zeta\Omega_k$\\
$\tau$ & dimensionless time, $c_rt/L$\\
$\Omega=\omega L/c_r$ & dimensionless circular frequency\\
$\kappa_r$, $\kappa_c$ & dimensionless wavenumbers, $\Omega$ and $\gamma\Omega$\\
$W_k$, $Y_k$, $m_k$, $\Gamma_k$ & mode shape, take-up amplitude, modal mass and participation factor of mode $k$\\
$H_A$, $H_B$; $\tilde m$ & receptances of strands A and B at the take-up; modal mass of a strand mode seen from its free end\\
$\tper{1}$, $\tauper{1}$ & fundamental period of the loop; $\tauper{1}=c_r\tper{1}/L$\\
$\tper{s}$, $\tauper{s}$ & fundamental period of a fixed--free strand\\
$\tauB$ & transit time of the strand downstream of the take-up, $1-\xi+\gamma$\\
$\betafive$ & take-up mass ratio at which a natural period lengthens by 5\,\%\\
\bottomrule
\end{tabular}
\end{table}


\subsection{Assumptions}\label{sec:assumptions}

The conveyor is idealised as a closed loop of belt driven by a single drive pulley
and tensioned by a single gravity take-up on the return strand. The following
assumptions are made.

\begin{enumerate}
  \item[(A1)] \emph{One-dimensional continuum.} The belt, the reduced (rotational)
  mass of the idlers and the part of the bulk material that moves with the belt are
  lumped into an equivalent inertial line density. The axial stiffness $EA$ is the
  same around the loop.
  \item[(A2)] \emph{Homogeneous strands.} The return and carry strands have constant
  inertial line densities $\mu_r$ and $\mu_c$ and therefore constant longitudinal
  wave speeds $c_j=\sqrt{EA/\mu_j}$, $j\in\{r,c\}$; both strands have the same path
  length $L$. The coupling of the bulk material to the belt is absorbed into
  $\mu_c$ through the coupling factor $\alpha$ of \citet{lodewijks2002}.
  \item[(A3)] \emph{Linear Kelvin--Voigt belt.} The dynamic tension is
  $T_d=EA\,(\varepsilon+t_v\,\dot\varepsilon)$, with retardation time $t_v$. The
  stiffening associated with belt sag at low tension \citep{ellis1987} and the
  dependence of the belt modulus on tension and temperature \citep{zarzycki2023}
  are neglected.
  \item[(A4)] \emph{Small belt speed.} The belt speed is two to three orders of
  magnitude below the wave speeds, so material and spatial coordinates coincide and
  transport terms are neglected.
  \item[(A5)] \emph{Point pulleys and rigid structure.} The contact between belt and
  pulley is reduced to a point. The rotational inertia and the bearing friction of
  the head, tail and bend pulleys are neglected, as are the elasticity and motion of
  the supporting structure.
  \item[(A6)] \emph{Drive.} The drive pulley does not slip and imposes a prescribed
  belt velocity $V(t)$. The no-slip condition is checked a posteriori
  (Section~\ref{sec:validity}); an imposed torque is treated as a limit of validity
  in Section~\ref{sec:torque}.
  \item[(A7)] \emph{Gravity take-up.} The take-up pulley is free to rotate and is
  carried by a carriage loaded by a counterweight. The translational masses of
  pulley, carriage and counterweight are condensed into a mass $\Mtu$ referred to
  the carriage displacement (Section~\ref{sec:static}). The rotational inertia of the
  take-up pulley, the elasticity of the rope, the inertia and friction of the
  sheaves, motion retarders and the end stops of the carriage are neglected.
  \item[(A8)] \emph{Linearity about the static state.} The conveyor starts from rest
  in static equilibrium, and the dynamic response is a linear increment on that
  state. The increment is admissible while the total belt tension stays positive
  (Section~\ref{sec:validity}).
\end{enumerate}

\subsection{Geometry and loop coordinate}\label{sec:geometry}

\begin{figure}[t]
  \centering
  \IfFileExists{figures/fig_loop.pdf}{%
    \includegraphics[width=\linewidth]{fig_loop}}{%
    \fbox{\parbox{0.9\linewidth}{\centering\vspace{3em}Loop schematic (run
    maps/fig\_loop.py).\vspace{3em}}}}
  \caption{Conveyor loop with the drive on the return strand. The geometric
  coordinate $\sigma$ runs from the head pulley along the return strand; the loop
  coordinate $s$ starts where the belt leaves the drive pulley ($s=0$) and ends where
  it returns to it ($s=2L$). The take-up at $s=\xi L$ divides the loop into strand~A,
  from the drive exit to the take-up, and strand~B, from the take-up by the tail, the
  carry strand and the head to the drive entry. Positions are illustrative.}
  \label{fig:loop}
\end{figure}

A geometric coordinate $\sigma\in[0,2L)$ is measured along the belt path in the
direction of travel, from the head pulley, so that the return strand occupies
$0\le\sigma<L$ and the carry strand $L\le\sigma<2L$; the tail pulley is at
$\sigma=L$ (\cref{fig:loop}). The drive pulley is at $\sigma_d$, anywhere on the
loop, and the take-up pulley at $\sigma_t\in(0,L)$, on the return strand. A head
drive corresponds to $\sigma_d=0$, a tail drive to $\sigma_d=L$ and an intermediate
drive to $0<\sigma_d<L$, with the take-up either downstream ($\sigma_t>\sigma_d$, on
the slack side) or upstream ($\sigma_t<\sigma_d$, on the tight side) of it.

The analysis uses the loop coordinate
\begin{equation}\label{eq:loopcoord}
  s=(\sigma-\sigma_d)\bmod 2L,\qquad s\in[0,2L],
\end{equation}
which starts where the belt leaves the drive pulley and ends where it returns to it;
$s=0$ and $s=2L$ are the two faces of the drive pulley. The take-up is at
$s_t=(\sigma_t-\sigma_d)\bmod 2L$, and $\xi=s_t/L$. In this coordinate the loop is an
open chain of homogeneous segments, separated by the head and tail pulleys, where the
line density changes, and by the take-up pulley. The take-up divides the chain into
strand~A ($0\le s\le s_t$) and strand~B ($s_t\le s\le 2L$). With a head drive,
strand~A is a piece of return belt of length $\xi L$, and strand~B holds the rest of
the return strand followed by the whole carry strand.

\subsection{Static state and dynamic increment}\label{sec:increment}

Let $\mw(s)$ be the line density that contributes to weight (belt and material,
without the reduced idler mass, so that $\mw\neq\mu$ in general), $\delta(s)$ the
local inclination of the belt path, positive uphill in the direction of travel, and
$\Tt$ the static belt tension at the take-up. At rest the static tension satisfies
\begin{equation}\label{eq:static}
  \frac{\dd T_0}{\dd s}=\mw(s)\,g\sin\delta(s),\qquad T_0(s_t)=\Tt,
\end{equation}
and the drive pulley, locked at rest, carries the difference $T_0(2L)-T_0(0)$. The
take-up holds $\Tt$ with the weight of its counterweight (Section~\ref{sec:static}).

During start-up the displacement of a belt particle is decomposed as
\begin{equation}\label{eq:decomposition}
  u(s,t)=u_0(s)+U(t)+w(s,t),\qquad U(t)=\int_0^tV(t')\,\dd t',
\end{equation}
where $u_0$ is the static displacement, $U$ the rigid-body motion imposed by the
drive and $w$ the elastic response. The rigid motion $U$ does not move the take-up,
because a uniform translation of the belt along its path leaves the length of belt
stored in the take-up unchanged. Balance of momentum of a belt element, after
subtraction of the static equilibrium~\eqref{eq:static}, gives in every segment
\begin{equation}\label{eq:pde}
  \mu(s)\,w_{tt}=\frac{\partial T_d}{\partial s}-\mu(s)\,a(t)-r(s)\,\varphi(t),
  \qquad T_d=EA\,(w_s+t_v\,w_{st}),
\end{equation}
with $a=\dot V$ the drive acceleration. The term $-\mu a$ is the inertial load of the
prescribed motion. The term $r(s)\varphi(t)$ represents the motion resistances
(idler rotation, indentation and flexure; force per unit length, opposing travel),
which are absent at rest and develop as the belt starts to move; $\varphi(t)$
describes their onset, with $\varphi\to1$ at steady speed. The base case is
$\varphi=V(t)/V_\infty$, with $V_\infty$ the final speed; a unit step serves as an
upper bound. A faithful representation would require resistances that depend on the
local belt speed, which is not attempted here.

Gravity does not appear in~\eqref{eq:pde}: it is balanced by the static state,
which includes the weight of the counterweight. Applying the weight as a load at
$t=0$ to a belt initially at rest and unstrained would excite a spurious transient of
the size of the static deflection. The simulations of \citet[Fig.~8]{lodewijks2013}
illustrate the superposition: during a drift stop of a \qty{1}{\kilo\metre} conveyor
whose path dips between head and tail, the dynamic drop of the carry-side tension in
the dip is the same, about \qty{17.3}{\kilo\newton}, for dips of
\qtyrange{0}{30}{\metre}, and the depth only shifts the static level. With a
\qty{40}{\metre} dip the tension drops to zero and the superposition fails, which is the slack limit of Section~\ref{sec:validity}. The
initial conditions are
\begin{equation}\label{eq:ic}
  w(s,0)=w_t(s,0)=0,\qquad y(0)=\dot y(0)=0,
\end{equation}
where $y$ is the displacement of the take-up pulley from its static position,
positive when the loop lengthens.

\subsection{Interface and boundary conditions}\label{sec:interfaces}

\paragraph{Drive pulley.} Without slip and with prescribed velocity, both faces of
the drive pulley follow $U(t)$; by~\eqref{eq:decomposition},
\begin{equation}\label{eq:bc_drive}
  w(0,t)=w(2L,t)=0.
\end{equation}

\paragraph{Head and tail pulleys.} By (A5) these pulleys transmit displacement and
tension unchanged,
\begin{equation}\label{eq:bc_ht}
  \jump{w}=0,\qquad \jump{T_d}=0,
\end{equation}
where $\jump{f}=f(s^+)-f(s^-)$. The same conditions hold at any other point where the
line density changes, such as the loading point.

\paragraph{Take-up pulley.} Consider first a pulley hung on the two belt branches
that wrap it by $180^\circ$ ($\nstr=2$). Because the pulley is free to rotate (A7),
the tension is the same on both sides,
\begin{equation}\label{eq:bc_tu_tension}
  T_d(s_t^-,t)=T_d(s_t^+,t)\equiv T_d(s_t,t).
\end{equation}
When the pulley moves by $y$, the belt stored around it grows by $2y$. In the loop
coordinate the take-up loop is collapsed to the point $s_t$, so the elongation of the
rest of the loop must grow by the same amount, $\int_0^{2L}w_s\,\dd s=2y$, the
integral being taken over the continuous parts. With~\eqref{eq:bc_drive},
\begin{equation}\label{eq:bc_tu_jump}
  \jump{w}_{s_t}=w(s_t^+,t)-w(s_t^-,t)=-2y(t):
\end{equation}
the belt upstream of the pulley is drawn forward into the take-up and the belt
downstream is drawn back. Finally, the two branches pull the pulley with $2T$; with
the static part balanced by the counterweight,
\begin{equation}\label{eq:bc_tu_dyn}
  \Mtu\,\ddot y=-2\,T_d(s_t,t).
\end{equation}
Equations~\eqref{eq:bc_tu_tension}--\eqref{eq:bc_tu_dyn} are the continuum
counterpart of the 2:1 take-up kinematics in the discrete models of
\citet[Eq.~7.55]{lodewijks1996} and \citet[Eqs.~4--5]{kulinowski2014}; in a continuum
model they were written by \citet[Eqs.~5--6]{song2012}, who then neglected the
take-up inertia. They follow from Hamilton's principle with kinetic energy
$\tfrac12\int\mu w_t^2\,\dd s+\tfrac12\Mtu\dot y^2$, strain energy
$\tfrac12\int EA\,w_s^2\,\dd s$ and the constraint~\eqref{eq:bc_tu_jump}: the
variation with respect to $w(s_t^\pm)$ gives~\eqref{eq:bc_tu_tension} and the
variation with respect to $y$ gives~\eqref{eq:bc_tu_dyn}. The static terms cancel
exactly, since the work of the static tension on the jump, $2\Tt y$, equals the work
of the take-up weight. The energy balance is
\begin{equation}\label{eq:energy}
  \frac{\dd}{\dd t}\left[\frac12\int_0^{2L}\!\left(\mu w_t^2+EA\,w_s^2\right)\dd s
  +\frac12\Mtu\dot y^2\right]
  =-\int_0^{2L}\!EA\,t_v\,w_{st}^2\,\dd s-\int_0^{2L}\!(\mu a+r\varphi)\,w_t\,\dd s .
\end{equation}

Eliminating $y$ between~\eqref{eq:bc_tu_jump} and~\eqref{eq:bc_tu_dyn} gives
$T_d(s_t)=(\Mtu/4)\,\jump{w_{tt}}_{s_t}$: the belt sees the take-up as a mass
$\Mtu/4$ in series, which resists only the relative acceleration of the two belt
ends and has no connection to the ground. This differs from a point mass embedded in
the belt (displacement continuous, tension jump $\jump{T_d}=\Mtu w_{tt}$), the
equivalent model of \citet[Fig.~2]{harrison1983} and the discrete model of
\citet{sakharwade2019}. The two models have opposite limits. As $\Mtu\to0$ the
embedded mass disappears, whereas the pulley enforces $T_d(s_t)=0$, a take-up of
constant force. As $\Mtu\to\infty$ the embedded mass becomes a fixed node, whereas
the pulley becomes a fixed pulley around which the belt runs freely and adds one slow
mode of the counterweight (Section~\ref{sec:properties}).

\subsection{Take-up hung on \texorpdfstring{$n$}{n} strands}\label{sec:nstrands}

Long conveyors often store belt in a multiple loop, with the carriage hung on
$\nstr$ belt strands ($\nstr$ even); $\nstr=4$ in the conveyor measured by
\citet{harrison1985b} and in the take-up described by \citet{nordell1997zisco}. With
carriage displacement $y_c$ and the same tension in all strands (free pulleys), the
conditions become $\jump{w}_{s_t}=-\nstr\,y_c$ and $\Mtu\ddot y_c=-\nstr\,T_d(s_t)$.
The equivalent displacement $y=(\nstr/2)\,y_c$ and the equivalent mass
\begin{equation}\label{eq:nmap}
  \Mtu_{\mathrm{eq}}=\frac{4\Mtu}{\nstr^2}
\end{equation}
recover~\eqref{eq:bc_tu_jump}--\eqref{eq:bc_tu_dyn} exactly and conserve the kinetic
energy, $\tfrac12\Mtu\dot y_c^2=\tfrac12\Mtu_{\mathrm{eq}}\dot y^2$. The static force on
the carriage is $\nstr\Tt$, the carriage travels $2/\nstr$ times the equivalent
displacement, and the length of belt stored, $2y$, does not depend on $\nstr$. The
number of take-up pulleys already enters the kinematics of the discrete model of
\citet{song2003}; here it only rescales the mass, and the rest of the paper is written
for the equivalent two-strand take-up. For the same take-up tension a multiple loop
reduces the equivalent mass by the factor $2/\nstr$ (Section~\ref{sec:static}), which
moves the take-up further towards the free-end regime of Section~\ref{sec:freeend}.

\subsection{Take-up mass and static tension}\label{sec:static}

The take-up mass is not a free parameter. Let $M_w$ be the counterweight mass, $M_c$
the mass of the take-up pulley and carriage, $\lambda$ the rigging ratio
(counterweight travel per unit carriage travel; $\lambda=1$ for a counterweight hung
directly on the carriage) and $\delta_c$ the inclination of the carriage track
($\delta_c=90^\circ$ for a vertical take-up). Virtual work and kinetic energy give
\begin{equation}\label{eq:reeving}
  \nstr\,\Tt=g\,(\lambda M_w+M_c\sin\delta_c),\qquad
  \Mtu=M_c+\lambda^2M_w=\frac{\lambda\,\nstr\,\Tt}{g}+(1-\lambda\sin\delta_c)\,M_c .
\end{equation}
With the equivalent mass~\eqref{eq:nmap} and the weight of the return strand
$\Wr=g\mu_rL$, the take-up mass ratio defined in Section~\ref{sec:dimensionless}
becomes
\begin{equation}\label{eq:beta_tension}
  \beta=\frac{4\Mtu}{\nstr^2\mu_rL}
  =\frac{4}{\nstr}\,\lambda\,\frac{\Tt}{\Wr}+\frac{4(1-\lambda\sin\delta_c)\,M_c}{\nstr^2\mu_rL}.
\end{equation}
The second term vanishes for a vertical take-up with a directly hung counterweight
and is otherwise small, since the carriage weighs much less than the counterweight.
Equation~\eqref{eq:beta_tension} states that the take-up mass follows from the
take-up tension: rigging and multiple loops change the mass for a given tension, but
a gravity take-up cannot combine a large tension with a small mass. The tension $\Tt$
is fixed by the running-state design: it is the smallest value that satisfies the
grip condition at the drive and the sag limits on both strands, plus whatever margin
the designer adds. It therefore depends weakly on the take-up position; for a
horizontal head-drive conveyor governed by grip, for example,
$\Tt=T_{2,\min}+r_r\,\xi L$, where $T_{2,\min}$ is the least admissible slack-side
tension and $r_r$ the running resistance per unit length of the return strand.
Section~\ref{sec:masscorrection} shows that the mass enters the dynamics only as a
correction, which is why the results are given in terms of $\xi$ and $\gamma$ with
$\beta$ as a correction; Section~\ref{sec:betaregime} evaluates $\Tt/\Wr$ for the
conveyors in the literature.

\subsection{Dimensionless formulation}\label{sec:dimensionless}

Lengths are scaled by $L$, times by the transit time of the return strand $L/c_r$,
accelerations by the peak drive acceleration $a_m$, displacements by $a_mL^2/c_r^2$
and tensions by the inertial force of the return strand $\mu_rLa_m$:
\begin{equation}\label{eq:scales}
  x=\frac{s}{L},\quad \tau=\frac{c_rt}{L},\quad
  \hat w=\frac{c_r^2}{a_mL^2}\,w,\quad \hat y=\frac{c_r^2}{a_mL^2}\,y,\quad
  \hat T=\frac{T_d}{\mu_rLa_m},\quad \hat a=\frac{a}{a_m},\quad
  \hat r=\frac{r}{\mu_ra_m}.
\end{equation}
The dimensionless parameters are
\begin{equation}\label{eq:params}
  \gamma=\frac{c_r}{c_c}=\sqrt{\frac{\mu_c}{\mu_r}},\qquad
  \beta=\frac{4\Mtu}{\nstr^2\mu_rL},\qquad
  \xi=\frac{s_t}{L},\qquad
  \hat\zeta=\frac{c_rt_v}{2L},\qquad
  \tau_a=\frac{c_rt_a}{L},
\end{equation}
together with the drive position $\sigma_d/L$ and the shape of the start-up profile,
whose acceleration time is $t_a$. For $\nstr=2$, $\beta=\Mtu/(\mu_rL)$. The relative
line density is $\hat\mu(x)=1$ on the return strand and $\gamma^2$ on the carry
strand. Since $EA=\mu_rc_r^2$, the problem~\eqref{eq:pde}--\eqref{eq:bc_tu_dyn}
becomes
\begin{gather}
  \hat\mu\,\hat w_{\tau\tau}=\frac{\partial\hat T}{\partial x}-\hat\mu\,\hat a(\tau)
  -\hat r(x)\,\varphi(\tau),\qquad \hat T=\hat w_x+2\hat\zeta\,\hat w_{x\tau},\label{eq:pde_nd}\\
  \hat w(0,\tau)=\hat w(2,\tau)=0,\qquad \jump{\hat w}_\xi=-2\hat y,\qquad
  \jump{\hat T}_\xi=0,\qquad \beta\,\hat y_{\tau\tau}=-2\,\hat T(\xi,\tau),\label{eq:bc_nd}
\end{gather}
with $\hat w$ and $\hat T$ continuous at the head and tail pulleys. In the elastic case
the dimensionless dynamic tension equals the dimensionless strain.
